% !TEX program = pdflatex
\documentclass[aspectratio=169]{beamer}

\usetheme{Madrid}
\usecolortheme{default}

\usepackage{amsmath, amssymb, bm, physics, mathtools}
\usepackage{braket}
\usepackage{bbm}
\usepackage{graphicx}
\usepackage{booktabs}
\usepackage{tikz}
\usepackage{quantikz}
\usepackage{hyperref}


% ----------------------------
% Convenience macros
% ----------------------------
\newcommand{\ketz}{\ket{0}}
\newcommand{\keto}{\ket{1}}
\newcommand{\ketzz}{\ket{00}}
\newcommand{\ketzo}{\ket{01}}
\newcommand{\ketoz}{\ket{10}}
\newcommand{\ketoo}{\ket{11}}

\newcommand{\Z}{\sigma_z}
\newcommand{\X}{\sigma_x}
\newcommand{\Y}{\sigma_y}
\newcommand{\I}{\mathbbm{1}}

\title[VQE for a Two-Qubit Hamiltonian]{How do we measure?}
\author{Morten Hjorth-Jensen}
\date{Spring 2026}

\begin{document}

% ============================================================
\begin{frame}
\titlepage
\end{frame}

% ============================================================
\begin{frame}{Two-qubit Hamiltonian as example case}
\begin{itemize}
  \item We are given a general two-qubit Hamiltonian \(H\), represented by a Hermitian \(4\times 4\) matrix.
  \item Goal: approximate the ground-state energy
  \[
    E_0 = \min_{\ket{\psi}\neq 0}\frac{\bra{\psi}H\ket{\psi}}{\braket{\psi}{\psi}}
    = \min_{\ket{\psi}, \ \|\psi\|=1} \bra{\psi}H\ket{\psi}.
  \]
  \item We will use the \textbf{Variational Quantum Eigensolver (VQE)}:
  \begin{enumerate}
    \item choose a parameterized state \(\ket{\psi(\bm{\theta})}\) prepared by a quantum circuit (ansatz),
    \item estimate \(E(\bm{\theta})=\bra{\psi(\bm{\theta})}H\ket{\psi(\bm{\theta})}\) from measurements,
    \item update parameters \(\bm{\theta}\) on a classical optimizer to (approximately) minimize \(E(\bm{\theta})\).
  \end{enumerate}
\end{itemize}
\end{frame}

% ============================================================
\begin{frame}{Measurement model: computational-basis (Pauli Z) readout}
Most devices naturally measure each qubit in the computational basis
\(\{\ket{0},\ket{1}\}\), i.e. eigenbasis of \(\Z\).

\medskip
A single two-qubit measurement shot yields a bitstring \(b_1 b_2\in\{00,01,10,11\}\).
We map outcomes to eigenvalues:
\[
\Z\ket{0}=+1\ket{0},\qquad \Z\ket{1}=-1\ket{1}.
\]
So define random variables per shot
\[
z_1 = \begin{cases}+1,& b_1=0\\-1,& b_1=1\end{cases},
\qquad
z_2 = \begin{cases}+1,& b_2=0\\-1,& b_2=1\end{cases}.
\]
Then the estimator for \(\expval{\Z\otimes\Z}\) is the sample mean of the product \(z_1 z_2\).
\textbf{Important:} Measuring \emph{both} qubits in the Z basis gives access to \(\Z\otimes \Z\) correlations in one shot.
\end{frame}

% ============================================================
\begin{frame}{Estimating \(\langle Z_1\rangle\), \(\langle Z_2\rangle\), \(\langle Z_1 Z_2\rangle\) from the same data}
Let \(p_{00},p_{01},p_{10},p_{11}\) be probabilities of outcomes \(\ket{00},\ket{01},\ket{10},\ket{11}\).
From Z-basis readout we have
\[
\expval{\Z\otimes \I}= (p_{00}+p_{01})-(p_{10}+p_{11}),
\]
\[
\expval{\I\otimes \Z}= (p_{00}+p_{10})-(p_{01}+p_{11}),
\]
\[
\expval{\Z\otimes \Z}= (p_{00}+p_{11})-(p_{01}+p_{10}).
\]

\medskip
\textbf{Practical estimator:} with \(N\) shots and counts \(n_{ij}\) (so \(p_{ij}\approx n_{ij}/N\)),
plug in \(p_{ij}=n_{ij}/N\) to obtain unbiased sample-mean estimators.
\end{frame}

% ============================================================
\begin{frame}{Why measuring both qubits is more convenient than only one}
Suppose you only measure qubit 1 in Z and discard qubit 2. You can estimate \(\expval{\Z\otimes\I}\), but:
\begin{itemize}
  \item You \textbf{cannot} directly estimate the two-qubit correlator \(\expval{\Z\otimes \Z}\) without joint outcomes.
  \item Many Hamiltonians include interaction terms like \(J\,\Z\otimes \Z\). Without joint readout, you would need additional circuitry or assumptions to reconstruct correlations.
  \item Joint measurement gives \(\expval{\Z\otimes\I}\), \(\expval{\I\otimes\Z}\), and \(\expval{\Z\otimes\Z}\) \textbf{simultaneously} from the same shot record \(\{b_1 b_2\}\).
\end{itemize}

\medskip
\textbf{Shot efficiency intuition:} if one measurement record yields multiple observables, you reduce total shots needed to estimate the energy to a given precision.

\medskip
\textbf{Physics intuition:} the ground-state energy of an interacting two-qubit system depends on \emph{correlations}; you must measure both subsystems to access them.
\end{frame}

% ============================================================
\begin{frame}{Measuring general Pauli strings using Z-basis readout}
Devices measure in Z, but VQE needs \(\expval{\X\otimes\X}\), \(\expval{\Y\otimes\Z}\), etc.

\medskip
\textbf{Strategy:} rotate the measurement basis so that measuring Z after a unitary \(U\) is equivalent to measuring another Pauli operator before \(U\).

\medskip
Single-qubit identities:
\[
H \Z H = \X, \qquad S^\dagger H \Z H S = \Y,
\]
where \(H\) is the Hadamard and \(S=\mathrm{diag}(1,i)\).

\medskip
Thus,
\begin{itemize}
  \item to measure \(\X\) on a qubit: apply \(H\) then measure in Z,
  \item to measure \(\Y\) on a qubit: apply \(S^\dagger\) then \(H\) then measure in Z,
  \item to measure \(\Z\): measure directly in Z.
\end{itemize}

For two-qubit strings like \(\X\otimes\Y\), apply the appropriate basis-change gates to each qubit and then measure both in Z.
\end{frame}

% ============================================================
\begin{frame}{Quantikz: measurement circuits for Pauli strings}
%\small
\textbf{Example 1:} measure \(\Z\otimes\Z\) (no basis change)
\[
\begin{quantikz}
\lstick{$\ket{\psi(\bm{\theta})}$} & \qw & \meter{} & \qw \\
\lstick{$\ket{\psi(\bm{\theta})}$} & \qw & \meter{} & \qw
\end{quantikz}
\]

\textbf{Example 2:} measure \(\X\otimes\Z\) (Hadamard on qubit 1)
\[
\begin{quantikz}
\lstick{$\ket{\psi(\bm{\theta})}$} & \gate{H} & \meter{} & \qw \\
\lstick{$\ket{\psi(\bm{\theta})}$} & \qw      & \meter{} & \qw
\end{quantikz}
\]
\end{frame}
\begin{frame}{Quantikz: measurement circuits for Pauli strings}
%\small

\textbf{Example 3:} measure \(\Y\otimes\Y\) (apply \(S^\dagger\) then \(H\) on each qubit)
\[
\begin{quantikz}
\lstick{$\ket{\psi(\bm{\theta})}$} & \gate{S^\dagger} & \gate{H} & \meter{} & \qw \\
\lstick{$\ket{\psi(\bm{\theta})}$} & \gate{S^\dagger} & \gate{H} & \meter{} & \qw
\end{quantikz}
\]
\end{frame}


% ============================================================
\begin{frame}{Setup and claim}
We consider a two-qubit state \(\rho\) and energy estimation for a Hamiltonian
\[
H = \sum_{\alpha,\beta\in\{0,x,y,z\}} h_{\alpha\beta}\, \sigma_\alpha\otimes\sigma_\beta,
\qquad \sigma_0=\I.
\]

\medskip
\textbf{Claim (informal).}
Measuring only the first qubit---even after arbitrary pre-rotations---cannot, in general, determine
two-qubit correlators such as \(\langle \Z\otimes\Z\rangle\), hence cannot determine \(\langle H\rangle\)
for generic interacting Hamiltonians.

\medskip
We will:
\begin{enumerate}
\item prove formally that any single-qubit measurement depends only on the reduced state \(\rho_1=\Tr_2\rho\),
\item exhibit distinct states with identical \(\rho_1\) but different \(\langle \Z\otimes\Z\rangle\),
\item give a worked VQE example where relying on single-qubit measurements provably fails.
\end{enumerate}
\end{frame}

% ============================================================
\begin{frame}{Formal fact 1: single-qubit statistics depend only on the reduced state}
Let \(\{M_m\}\) be a POVM (s a measure whose values are positive semi-definite operators on a Hilbert space) performed \emph{only} on qubit 1. Operationally this corresponds to
measurement operators \(\{M_m\otimes \I\}\) on the two-qubit system.

\medskip
The probability of outcome \(m\) is
\[
p(m)=\Tr\!\big[(M_m\otimes \I)\rho\big].
\]
Using the defining property of the partial trace,
\[
\Tr\!\big[(A\otimes \I)\rho\big] = \Tr\!\big[A\,\Tr_2(\rho)\big],
\]
we obtain
\[
p(m)=\Tr\!\big[M_m\,\rho_1\big],\qquad \rho_1:=\Tr_2(\rho).
\]

\medskip
\textbf{Conclusion.} \emph{All} statistics from measurements that act only on qubit 1 depend only on \(\rho_1\).
They are completely insensitive to correlations that live in \(\rho\) but vanish under \(\Tr_2\).
\end{frame}


% ============================================================
\begin{frame}{Corollary: no single-qubit readout can determine two-qubit correlators}
A two-qubit correlator (e.g. \(\langle \Z\otimes\Z\rangle\)) is
\[
\langle \Z\otimes\Z\rangle_\rho=\Tr\!\big[(\Z\otimes\Z)\rho\big].
\]

From the previous slide, any single-qubit protocol (measurements on qubit 1 only) can at best determine \(\rho_1\).
Therefore, if there exist \(\rho,\rho'\) such that
\[
\Tr_2(\rho)=\Tr_2(\rho') \quad \text{but} \quad
\Tr\!\big[(\Z\otimes\Z)\rho\big]\neq \Tr\!\big[(\Z\otimes\Z)\rho'\big],
\]
then \(\langle \Z\otimes\Z\rangle\) cannot be inferred from single-qubit measurement data.

\medskip
Next slide: we construct such a pair explicitly (pure states, maximally entangled).
\end{frame}

% ============================================================
\begin{frame}{Exhibit: same reduced state, different \(\langle Z\otimes Z\rangle\)}
Consider the Bell states
\[
\ket{\Phi^+}=\frac{\ket{00}+\ket{11}}{\sqrt{2}},\qquad
\ket{\Psi^+}=\frac{\ket{01}+\ket{10}}{\sqrt{2}}.
\]
Let \(\rho_{\Phi^+}=\ket{\Phi^+}\!\bra{\Phi^+}\) and \(\rho_{\Psi^+}=\ket{\Psi^+}\!\bra{\Psi^+}\).

\medskip
\textbf{Compute reduced states.}
For either state,
\[
\rho_1=\Tr_2(\rho)=\frac{\I}{2}.
\]
Hence \emph{every} measurement on qubit 1 alone yields identical statistics for \(\ket{\Phi^+}\) and \(\ket{\Psi^+}\).

\medskip
\textbf{Compute the correlation.}
Using \(\Z\ket{0}=+\ket{0}\), \(\Z\ket{1}=-\ket{1}\),
\[
(\Z\otimes\Z)\ket{00}=+\,\ket{00},\quad (\Z\otimes\Z)\ket{11}=+\,\ket{11},
\]
\[
(\Z\otimes\Z)\ket{01}=-\,\ket{01},\quad (\Z\otimes\Z)\ket{10}=-\,\ket{10}.
\]
Therefore,
\[
\langle \Z\otimes\Z\rangle_{\Phi^+}=+1,\qquad
\langle \Z\otimes\Z\rangle_{\Psi^+}=-1.
\]

\medskip
\textbf{Key consequence.} Single-qubit measurements cannot distinguish \(\ket{\Phi^+}\) from \(\ket{\Psi^+}\),
yet the interaction observable \(\Z\otimes\Z\) takes opposite values.
\end{frame}

% ============================================================
\begin{frame}{Formal theorem: impossibility of recovering generic energies from one-qubit measurements}
\textbf{Theorem.}
Let \(H\) be any Hamiltonian that contains at least one nonzero two-qubit term, e.g. \(h_{zz}\,\Z\otimes\Z\) with \(h_{zz}\neq 0\).
Then there exist states \(\rho,\rho'\) such that:
\[
\Tr_2(\rho)=\Tr_2(\rho')\quad\text{but}\quad \Tr(H\rho)\neq \Tr(H\rho').
\]
Hence the energy \(\langle H\rangle\) cannot be determined from measurements on qubit 1 alone.

\medskip
\textbf{Proof (constructive).}
Take \(\rho=\rho_{\Phi^+}\), \(\rho'=\rho_{\Psi^+}\). Then \(\Tr_2(\rho)=\Tr_2(\rho')=\I/2\),
but for the \(h_{zz}\) term,
\[
\Tr\!\big[(\Z\otimes\Z)\rho\big]=+1,\qquad \Tr\!\big[(\Z\otimes\Z)\rho'\big]=-1.
\]
Therefore
\[
\Tr(H\rho)-\Tr(H\rho') \supset h_{zz}\big(+1-(-1)\big)=2h_{zz}\neq 0,
\]
so \(\Tr(H\rho)\neq \Tr(H\rho')\), while all one-qubit statistics on qubit 1 are identical.
\(\square\)
\end{frame}

% ============================================================
\begin{frame}{Addressing the ``pre-rotation'' idea: why it doesn't save single-qubit readout}
One may attempt: apply a two-qubit unitary \(U\), then measure only qubit 1 in the Z basis.
This measures the observable
\[
A := U^\dagger(\Z\otimes\I)U
\]
on the original state.

\medskip
\textbf{Two key limitations:}
\begin{itemize}
\item \textbf{Single observable per circuit.} Measuring only qubit 1 gives information about expectation values of \emph{one} operator \(A\) per measurement setting.
VQE generally needs many Pauli strings \(\{\sigma_\alpha\otimes\sigma_\beta\}\).
\item \textbf{No universal compression.} There is no fixed unitary \(U\) such that
\(U^\dagger(\Z\otimes\I)U\) equals \(\sigma_\alpha\otimes\sigma_\beta\) for \emph{all} required terms simultaneously.
You would need a different \(U\) per term anyway, eliminating any advantage over simply measuring both qubits.
\end{itemize}

\medskip
\textbf{Conceptual punchline:} two-qubit Hamiltonians depend on correlations; correlations are not encoded in a single reduced state.
\end{frame}

% ============================================================
\begin{frame}{Worked VQE failure example: Hamiltonian requiring correlations}
Consider the interacting Hamiltonian
\[
H = -\,\Z\otimes\Z.
\]
Its eigenvalues are \(\{-1,-1,+1,+1\}\) with ground energy \(E_0=-1\).
Any ground state must satisfy \(\langle \Z\otimes\Z\rangle=+1\)
(e.g. \(\ket{00}\), \(\ket{11}\), or any superposition in the \(+1\) eigenspace).

\medskip
\textbf{Energy functional:}
\[
E(\rho)=\Tr(H\rho) = -\,\langle \Z\otimes\Z\rangle_\rho.
\]

\medskip
Now compare the two Bell states from earlier:
\[
E(\Phi^+) = -\,(+1)=-1,\qquad E(\Psi^+)=-\,(-1)=+1.
\]
They differ by \(\Delta E=2\), but have identical \(\rho_1=\I/2\).

\medskip
Therefore any VQE evaluation strategy that uses only qubit-1 measurement data \emph{cannot even tell whether the energy is \(-1\) or \(+1\)}.
\end{frame}



% ============================================================
\begin{frame}{VQE ansatz family demonstrating the failure concretely}
Take a simple two-qubit ansatz capable of preparing both Bell states:
\[
\ket{\psi(\theta)} = U(\theta)\ket{00},
\quad
U(\theta) := \mathrm{CNOT}_{1\to 2}\,(R_y(\theta)\otimes \I).
\]

One checks:
\[
\theta=\frac{\pi}{2}\ \Rightarrow\ \ket{\psi(\theta)}=\ket{\Phi^+},
\qquad
\theta=\frac{\pi}{2}\ \text{with an extra }X\text{ on qubit 2}\Rightarrow \ket{\Psi^+}.
\]
\medskip
But for both \(\ket{\Phi^+}\) and \(\ket{\Psi^+}\),
\[
\rho_1=\Tr_2(\ket{\psi}\!\bra{\psi})=\frac{\I}{2}
\quad\Rightarrow\quad
\langle \Z\otimes\I\rangle = \Tr(\Z\rho_1)=0,
\]
and similarly any single-qubit observable on qubit 1 has the same expectation (because \(\rho_1=\I/2\)).

\medskip
\textbf{Thus:} if your VQE ``measurement layer'' reads only qubit 1, the cost function becomes \emph{constant} across this family and provides no gradient/no signal for optimization.
\end{frame}

% ============================================================
\begin{frame}{What measurement on both qubits provides (and why it fixes VQE)}
If you measure \emph{both} qubits in the computational basis, each shot yields \(b_1b_2\in\{00,01,10,11\}\).
Map bits to \(\pm 1\) eigenvalues:
\[
z_1=(-1)^{b_1},\qquad z_2=(-1)^{b_2}.
\]
Then a single shot provides:
\[
z_1 \Rightarrow \Z\otimes\I,\qquad
z_2 \Rightarrow \I\otimes\Z,\qquad
z_1z_2 \Rightarrow \Z\otimes\Z.
\]

\medskip
So the estimator for \(\langle \Z\otimes\Z\rangle\) is
\[
\widehat{\langle \Z\otimes\Z\rangle}=\frac{1}{N}\sum_{s=1}^N z_1^{(s)}z_2^{(s)},
\]
and the energy estimate is
\[
\hat{E} = -\,\widehat{\langle \Z\otimes\Z\rangle}.
\]

\medskip
\textbf{Efficiency point:} the same shot record gives multiple observables at once; measuring only one qubit throws away correlation information and wastes shots.
\end{frame}

% ============================================================
\begin{frame}{Generalization: why this is not a special case}
A generic two-qubit Hamiltonian has the Pauli expansion
\[
H=\sum_{\alpha,\beta} h_{\alpha\beta}\,\sigma_\alpha\otimes\sigma_\beta,
\]
with up to 15 nontrivial coefficients \((\alpha,\beta)\neq (0,0)\).

\medskip
\textbf{If \(H\) contains any interaction term} (e.g. \(\Z\otimes\Z\), \(\X\otimes\X\), \(\Y\otimes\Y\), \(\X\otimes\Z\), etc.),
then the energy depends on correlators \(\langle \sigma_\alpha\otimes\sigma_\beta\rangle\).

\medskip
But measuring only qubit 1 gives access only to \(\langle \sigma_\alpha\otimes \I\rangle\), i.e. local expectations.

\medskip
\textbf{Formal obstruction:} There exist inequivalent two-qubit states sharing the same \(\rho_1\),
so no function of \(\rho_1\) alone can reproduce all correlators.
Thus single-qubit measurement cannot generally yield the correct energy landscape for VQE.
\end{frame}

% ============================================================
\begin{frame}{More examples of measurement circuits}
Modern devices measure all qubits in the Z basis; other Pauli strings are measured via pre-rotations.
\textbf{(i) \(X\otimes X\):} apply \(H\) on both, then Z readout
\[
\begin{quantikz}
\lstick{$\ket{\psi(\bm{\theta})}$} & \gate{H} & \meter{} & \qw \\
\lstick{$\ket{\psi(\bm{\theta})}$} & \gate{H} & \meter{} & \qw
\end{quantikz}
\]

\textbf{(ii) \(Y\otimes Y\):} apply \(S^\dagger\) then \(H\) on both, then Z readout
\[
\begin{quantikz}
\lstick{$\ket{\psi(\bm{\theta})}$} & \gate{S^\dagger} & \gate{H} & \meter{} & \qw \\
\lstick{$\ket{\psi(\bm{\theta})}$} & \gate{S^\dagger} & \gate{H} & \meter{} & \qw
\end{quantikz}
\]
\end{frame}


% ============================================================
\begin{frame}{Why joint Z readout is sufficient for energy estimation}
Because each Pauli string measurement is reduced to:
\begin{enumerate}
  \item local basis-change gates \(U_k = U_k^{(1)}\otimes U_k^{(2)}\),
  \item joint Z readout producing bitstrings \(b_1 b_2\),
\end{enumerate}
we have
\[
\expval{P_k}_{\psi}
=
\expval{(\Z^{p_1}\otimes \Z^{p_2})}_{U_k\psi}
\quad\text{with}\quad p_i\in\{0,1\},
\]
where the mapping depends on whether the rotated measurement corresponds to \(\X,\Y,\Z\).

\textbf{Operationally:} you always measure in the computational basis, but you interpret the bits as \(\pm 1\) eigenvalues (after the appropriate pre-rotations).

\textbf{And:} joint readout gives you correlators (e.g. \(Z_1 Z_2\)) needed for interaction energies.
\end{frame}


% ============================================================
\begin{frame}{(Optional) Operator grouping and measurement efficiency}
Many Pauli strings \(\{P_k\}\) commute and can be measured in the same basis.
For two qubits, an important commuting set is:
\[
\{\Z\otimes\I,\ \I\otimes\Z,\ \Z\otimes\Z\}
\]
which all share the Z basis and are estimated from the \emph{same} Z-basis shots.

\medskip
Similarly, \(\{\X\otimes\X,\ \X\otimes\I,\ \I\otimes\X\}\) share the X basis (apply \(H\) on relevant qubits),
and analogous for Y basis.

\medskip
\textbf{Benefit:} fewer distinct measurement circuits \(\Rightarrow\) fewer total shots for a target energy precision.
\end{frame}

% ============================================================
\begin{frame}{Summary}
\begin{itemize}
  \item Any two-qubit Hermitian \(4\times 4\) Hamiltonian admits a Pauli expansion \(H=\sum h_{\alpha\beta}\sigma_\alpha\otimes\sigma_\beta\).
  \item VQE minimizes \(E(\bm{\theta})=\bra{\psi(\bm{\theta})}H\ket{\psi(\bm{\theta})}\) using the variational principle.
  \item Energy estimation reduces to estimating expectation values of Pauli strings.
  \item All measurements can be performed with \textbf{computational-basis (Z)} readout by inserting local basis-change gates.
  \item Measuring \textbf{both qubits} is more convenient because it yields \(\Z\otimes\Z\) (and other two-qubit correlators) directly and improves shot efficiency by extracting multiple observables from one record.
\end{itemize}
\end{frame}

% ============================================================
\begin{frame}{More Summary (what you should remember)}
\begin{itemize}
\item \textbf{Formal fact:} any measurement acting only on qubit 1 depends only on the reduced state \(\rho_1=\Tr_2\rho\).
\item \textbf{Correlations are invisible locally:} there exist states with the same \(\rho_1\) but different \(\langle \sigma_\alpha\otimes\sigma_\beta\rangle\).
\item \textbf{VQE requires correlators} because interacting Hamiltonians include two-qubit terms at least.
\item \textbf{Practical takeaway:} measuring both qubits in the Z basis yields local expectations and correlators from the same bitstrings, enabling correct and shot-efficient energy estimation.
\end{itemize}
\end{frame}

\end{document}



\usepackage{pgf}
\usepackage{graphicx}
\usepackage{epsfig}
\usepackage{relsize}

\usepackage{fancybox}  % make sure fancybox is loaded before fancyvrb

\usepackage{fancyvrb}
%\usepackage{minted} % requires pygments and latex -shell-escape filename
%\usepackage{anslistings}
%\usepackage{listingsutf8}

\usepackage{amsmath,amssymb,bm}
%\usepackage[latin1]{inputenc}
\usepackage[T1]{fontenc}
\usepackage[utf8]{inputenc}
\usepackage{colortbl}
\usepackage[english]{babel}
\usepackage{tikz}
\usepackage{framed}
% Use some nice templates
\beamertemplatetransparentcovereddynamic

% --- begin table of contents based on sections ---
% Delete this, if you do not want the table of contents to pop up at
% the beginning of each section:
% (Only section headings can enter the table of contents in Beamer
% slides generated from DocOnce source, while subsections are used
% for the title in ordinary slides.)
\AtBeginSection[]
{
  \begin{frame}<beamer>[plain]
  \frametitle{}
  %\frametitle{Outline}
  \tableofcontents[currentsection]
  \end{frame}
}


\title{Quantum Computing, Quantum Machine Learning and Quantum Information Theories}

% ----------------- author(s) -------------------------

\author{Morten Hjorth-Jensen\inst{1}}
\institute{Department of Physics, University of Oslo, Norway\inst{1}}
% ----------------- end author(s) -------------------------

\date{February 25, 2026
% <optional titlepage figure>
\ \\ 
{\tiny \copyright\ 1999-2026, Morten Hjorth-Jensen. Released under CC Attribution-NonCommercial 4.0 license}
}

\begin{frame}[plain,fragile]
\titlepage
\end{frame}

\begin{frame}[plain,fragile]
\frametitle{Plans for the week of February 23-27}

\begin{block}{}
\begin{itemize}
\item Repetition from last week on the VQE with applications to a one-qubit problem

\item Setting up calculations with the Variational Quantum Eigensolver (VQE) for one- and two-qubits Hamiltonians

\item Deriving the expressions for the gradients and setting up measurements in the computational basis 
% * Video of lecture at \href{{https://youtu.be/}}{\nolinkurl{https://youtu.be/}}
% * Whiteboard notes at \href{{https://github.com/CompPhysics/QuantumComputingMachineLearning/blob/gh-pages/doc/HandWrittenNotes/2026/Notesweek6.pdf}}{\nolinkurl{https://github.com/CompPhysics/QuantumComputingMachineLearning/blob/gh-pages/doc/HandWrittenNotes/2026/Notesweek6.pdf}}
\end{itemize}

\noindent
\end{block}
\end{frame}

\begin{frame}[plain,fragile]
\frametitle{Reading suggestions}

\begin{block}{}
\begin{itemize}
\item Hundt's text section 6.11 on the VQE

\item Jupyter-notebook on VQE and single-qubit problem at \href{{https://github.com/CompPhysics/QuantumComputingMachineLearning/blob/gh-pages/doc/pub/week6/ipynb/single_qubit_vqe_modified.ipynb}}{\nolinkurl{https://github.com/CompPhysics/QuantumComputingMachineLearning/blob/gh-pages/doc/pub/week6/ipynb/single_qubit_vqe_modified.ipynb}}

\item Jupyter-notebook on VQE and two-qubit problem at \href{{https://github.com/CompPhysics/QuantumComputingMachineLearning/blob/gh-pages/doc/pub/week6/ipynb/two_qubit_vqe_modified.ipynb}}{\nolinkurl{https://github.com/CompPhysics/QuantumComputingMachineLearning/blob/gh-pages/doc/pub/week6/ipynb/two_qubit_vqe_modified.ipynb}}
\end{itemize}

\noindent
\end{block}
\end{frame}

\begin{frame}[plain,fragile]
\frametitle{VQE overview}

\vspace{6mm}

% inline figure
\centerline{\includegraphics[width=1.0\linewidth]{figures/vqe.png}}

\vspace{6mm}
\end{frame}

\begin{frame}[plain,fragile]
\frametitle{One-qubit Hamiltonian}

Last week we simply computed the expectations values of the various
Pauli matrices by evaluating expressions like
\[
\langle \psi(\theta)\vert \bm{X} \vert \psi(\theta)\rangle,
\]
for the Pauli-$\bm{X}$ matrix and similarly for the other Pauli matrices and
the $2\times 2$ identity matrix $\bm{I}$.  This is however not the way
we end up evaluating such expectation values on an actual quantum
computer. We have to transform the set of measurements to be performed to the computational basis we have chosen.
\end{frame}

\begin{frame}[plain,fragile]
\frametitle{Our computational basis}

The computational basis we have chosen is given by the eigenvectors of
the Pauli-$\bm{Z}$ matrix, namely $\vert 0 \rangle$ and $\vert
1\rangle$.

For the actual implementation of these measurements we need three different circuits where
\begin{itemize}
\item $\bm{Z}$-measurement: Direct measurement.

\item $\bm{X}$-measurement: Apply a Hadamard gate $\bm{H}$, then measure $\bm{Z}$.

\item $\bm{Y}$-measurement: Apply $\bm{S}^{\dagger}\bm{H}$ gates, then measure $\bm{Z}$.
\end{itemize}

\noindent
\end{frame}

\begin{frame}[plain,fragile]
\frametitle{Shot-Based Sampling}

\begin{block}{Finite measurement shots introduce statistical noise }
\begin{itemize}
\item Each measurement yields 0 or 1 with probabilities $|\alpha|^2$ and $|\beta|^2$

\item Expectation values estimated from shot counts

\item Error scales as $\mathcal{O}(1/\sqrt{N_{\text{shots}}})$
\end{itemize}

\noindent
\end{block} 
The codes we discussed last week did not include a shot-based sampling.
\end{frame}

\begin{frame}[plain,fragile]
\frametitle{Parameter-shift rule for evaluating gradients}

In our discussions last week we used a quasi-Newton methods to find
the roots of the gradients, meaning that we evaluated the gradients by
finite difference methods.

We will now derive a more efficient way of evaluating the derivatives
of the energy with respect to the parameters of our circuit (the
rotation angles).
\end{frame}

\begin{frame}[plain,fragile]
\frametitle{Gradient descent and calculations of gradients}

In order to optimize the VQE ansatz, we need to compute derivatives
with respect to the variational parameters.  Here we develop first a
simpler approach tailored to the one-qubit case. For this particular
case, we have defined an ansatz in terms of the Pauli rotation
matrices.
\end{frame}

\begin{frame}[plain,fragile]
\frametitle{Setting up gradients}

These define an arbitrary one-qubit state on the Bloch
sphere through the expression

\[
\vert\psi\rangle = \vert \psi(\theta,\phi)\rangle =R_y(\phi)R_x(\theta)\vert 0 \rangle.
\]
Each of these rotation matrices can be written in a more general form as
\[
R_{i}(\gamma)=\exp{-(\imath\frac{\gamma}{2}\sigma_i)}=\cos{(\frac{\gamma}{2})}\bm{I}-\imath\sin{(\frac{\gamma}{2})}\bm{\sigma}_i,
\]
where $\sigma_i$ is one of the Pauli matrices $\sigma_{x,y,z}$.
\end{frame}

\begin{frame}[plain,fragile]
\frametitle{Derivatives}

It is easy to see that the derivative with respect to $\gamma$ is
\[
\frac{\partial R_{i}(\gamma)}{\partial \gamma}=-\frac{\gamma}{2}\bm{\sigma}_i R_{i}(\gamma).
\]
\end{frame}

\begin{frame}[plain,fragile]
\frametitle{Derivatives of the expectation value of the Hamiltonian}

We can now calculate the derivative of the expectation value of the
Hamiltonian in terms of the angles $\theta$ and $\phi$. We have two
derivatives 
\[
\frac{\partial}{\partial \theta}\left[\langle \psi(\theta,\phi) \vert \bm{H}\vert \psi(\theta,\phi)\rangle\right]=\frac{\partial}{\partial \theta}\left[\langle\bm{H}(\theta,\phi)\rangle\right]=\langle \psi(\theta,\phi) \vert \bm{H}(-\frac{\imath}{2}\bm{\sigma}_x\vert \psi(\theta,\phi)\rangle+\hspace{0.1cm}\mathrm{h.c},
\]
and
\[
\frac{\partial }{\partial \phi}\left[\langle \psi(\theta,\phi) \vert \bm{H}\vert \psi(\theta,\phi)\rangle\right]=\frac{\partial}{\partial \phi}\left[\langle\bm{H}(\theta,\phi)\rangle\right]=\langle \psi(\theta,\phi) \vert \bm{H}(-\frac{\imath}{2}\bm{\sigma}_y\vert \psi(\theta,\phi)\rangle+\hspace{0.1cm}\mathrm{h.c}. 
\]
\end{frame}

\begin{frame}[plain,fragile]
\frametitle{Two addtional expectation values}

This means that we have to calculate two additional expectation values
in addition to the expectation value of the Hamiltonian itself.  If we
stay with an ansatz for the single qubit states given by the above
rotation operators, we can, following for example \href{{https://journals.aps.org/pra/abstract/10.1103/PhysRevA.99.032331}}{the article by
Maria Schuld et
al},
show that the derivative of the expectation value of the Hamiltonian
can be written as (we focus only on a given angle $\phi$)

\[
\frac{\partial}{\partial \phi}\left[\langle\bm{H}(\phi)\rangle\right]=\frac{1}{2}\left[\langle\bm{H}(\phi+\frac{\pi}{2})\rangle-\langle\bm{H}(\phi-\frac{\pi}{2})\rangle\right].
\]
\end{frame}

\begin{frame}[plain,fragile]
\frametitle{Rotations again and again}

To see this, consider again the definition of the rotation operators.
We can write these operators as

\[
R_i(\phi)=\exp{-\imath(\phi \bm{\sigma}_i)},
\]
with $\bm{sigma}_i$, with $\bm{\sigma}_i$ being any of the Pauli
matrices $\bm{X}$, $\bm{Y}$ and $\bm{Z}$. The latter can be generalized to other
unitary matrices as well.
The derivative with respect to $\phi$ gives
\[
\frac{\partial R_i(\phi)}{\partial \phi}=-\frac{\imath}{2} \bm{\sigma}_i\exp{-\imath(\phi \bm{\sigma}_i)}=-\frac{\imath}{2} \bm{\sigma} R_i(\phi).
\]
\end{frame}

\begin{frame}[plain,fragile]
\frametitle{Bloch sphere math}

Our ansatz for a general one-qubit state on the Bloch sphere contains the product of a rotation around the $x$-axis and the $y$-axis. In the derivation here we focus only on one angle however. Our ansatz is then given by

\[
\vert \psi \rangle = R_i(\phi)\vert 0 \rangle,
\]
and the expectation value of our Hamiltonian is

\[
\langle \psi \vert \mathcal{H}\vert \psi \rangle = \langle 0 \vert R_i(\phi)^{\dagger} \mathcal{H}R_i(\phi)\vert 0\rangle. 
\]
\end{frame}

\begin{frame}[plain,fragile]
\frametitle{Derivatives}

Our derivative with respect to the angle $\phi$ has a similar structure, that is
\[
\frac{\partial }{\partial \phi}\left[\langle \psi(\theta,\phi) \vert \bm{H}\vert \psi(\theta,\phi)\rangle\right]=\langle \psi(\theta,\phi) \vert \bm{H}(-\frac{\imath}{2}\bm{\sigma}_y\vert \psi(\theta,\phi)\rangle+\hspace{0.1cm}\mathrm{h.c}. 
\]
\end{frame}

\begin{frame}[plain,fragile]
\frametitle{Rewriting}

In order to rewrite the equation of the derivative,
the following relation is useful
\[
\langle \psi \vert \hat{A}^{\dagger}\hat{B}\hat{C}\vert \psi \rangle = \frac{1}{2}\left[
\langle \psi \vert (\hat{A}+\hat{C})^{\dagger}\hat{B}(\bm{A}+\hat{C})\vert \psi \rangle-\langle \psi \vert (\hat{A}-\hat{C})^{\dagger}\hat{B}(\bm{A}-\hat{C})\vert \psi \rangle\right],
\]

where $\hat{A}$, $\hat{B}$ and $\hat{C}$ are arbitrary hermitian
operators.
\end{frame}

\begin{frame}[plain,fragile]
\frametitle{Final manipulations}

If we identify these operators as $\hat{A}=\bm{I}$, with
$\bm{I}$ being the unit operator, $\hat{B}=\mathcal{H}$ our Hamiltonian,
and $\hat{C}=-\imath \bm{\sigma}_i/2$, we obtain the following
expression for the expectation value of the derivative (excluding the hermitian conjugate)

\[
\langle \psi \vert \bm{I}^{\dagger}\mathcal{H}(-\frac{\imath}{2}\bm{\sigma}_i\vert \psi \rangle = \frac{1}{2}\left[
\langle \psi \vert (\bm{I}-\frac{\imath}{2} \bm{\sigma}_i)^{\dagger}\mathcal{H}(\bm{I}-\frac{\imath}{2} \bm{\sigma}_i)\vert \psi \rangle-\langle \psi \vert (\bm{I}+\frac{\imath}{2} \bm{\sigma}_i)^{\dagger}\mathcal{H}(\bm{I}+\frac{\imath}{2} \bm{\sigma}_i)\vert \psi \rangle\right].
\]
\end{frame}

\begin{frame}[plain,fragile]
\frametitle{The expressions to implement}

If we then use that the rotation matrices can be rewritten as
\[
R_{i}(\phi)=\exp{-(\imath\frac{\phi}{2}\sigma_i)}=\cos{(\frac{\phi}{2})}\bm{I}-\imath\sin{(\frac{\phi}{2})}\bm{\sigma}_i,
\]
we see that if we set the angle to $\phi=\pi/2$, we have 
\[
R_{i}(\frac{\pi}{2})=\cos{(\frac{\pi}{4})}\bm{I}-\imath\sin{(\frac{\pi}{4})}\bm{\sigma}_i=\frac{1}{\sqrt{2}}\left(\bm{I}-\frac{\imath}{2} \bm{\sigma}_i\right).
\]
\end{frame}

\begin{frame}[plain,fragile]
\frametitle{Final expression}

This means that we can write
\[
\langle \psi \vert \bm{I}^{\dagger}\mathcal{H}(-\frac{\imath}{2}\bm{\sigma}_i\vert \psi \rangle = \frac{1}{2}\left[
\langle \psi \vert R_i(\frac{\pi}{2})^{\dagger}\mathcal{H}R_i(\frac{\pi}{2})\vert \psi \rangle-\langle \psi \vert R_i(-\frac{\pi}{2})^{\dagger}\mathcal{H}R_i(-\frac{\pi}{2})^{\dagger}\vert \psi \rangle\right]=\frac{1}{2}(\langle\mathcal{H}(\phi+\frac{\pi}{2})\rangle-\langle\mathcal{H}(\phi-\frac{\pi}{2})\rangle).
\]
\end{frame}

\begin{frame}[plain,fragile]
\frametitle{Basics of gradient descent and stochastic gradient descent}

In order to implement the above equations, we need to remind the
reader about basic elements of various optimization approaches. Our
main focus here will be various gradient descent approaches and quasi-Newton methods like Broyden's algorithm and variations thereof.

This material is covered by the lectures from \href{{https://github.com/CompPhysics/ComputationalPhysics2/blob/gh-pages/doc/pub/week6/ipynb/week6.ipynb}}{FYS4411 on gradient optimization}
\end{frame}

\begin{frame}[plain,fragile]
\frametitle{Two-qubit Hamiltonian}

Here we discuss  how to rewrite the two-qubit Hamiltonian defined by the following Hamiltonian matrix (project 1)
\[
\mathcal{H}=\begin{bmatrix} \epsilon_{1}+V_z & 0 & 0 & V_x \\
                       0  & \epsilon_{2}-V_z & V_x & 0 \\
		       0 & H_x & \epsilon_{3}-V_z & 0 \\
		       H_x & 0 & 0 & \epsilon_{4} +V_z \end{bmatrix}.
\] 

This Hamiltonian can be rewritten in terms of various one-qubit matrices.
\end{frame}

\begin{frame}[plain,fragile]
\frametitle{Definitions}

We define
\[
\epsilon_{II}=\frac{\epsilon_{1}+\epsilon_{2}+\epsilon_{3}+\epsilon_{4}}{4},
\]
\[
\epsilon_{ZI}=\frac{\epsilon_{1}+\epsilon_{2}-\epsilon_{3}-\epsilon_{4}}{4},
\]
\[
\epsilon_{IZ}=\frac{\epsilon_{1}-\epsilon_{2}+\epsilon_{3}-\epsilon_{4}}{4},
\]
\[
\epsilon_{ZZ}=\frac{\epsilon_{1}-\epsilon_{2}-\epsilon_{3}+\epsilon_{4}}{4}.
\]
\end{frame}

\begin{frame}[plain,fragile]
\frametitle{The Hamiltonian in terms of Pauli-$\bm{X}$ and Pauli-$\bm{Z}$ matrices}

With these definitions we can rewrite our two-qubit Hamiltonian as
\[
\mathcal{H}=\mathcal{H}_0+\mathcal{H}_I
\]
with
\[
\mathcal{H}_0=\epsilon_{II}\bm{I}\otimes\bm{I}+\epsilon_{ZI}\bm{Z}\otimes\bm{I}+\epsilon_{IZ}\bm{I}\otimes\bm{Z}+\epsilon_{ZZ}\bm{Z}\otimes\bm{Z},
\]
and
\[
\mathcal{H}_I=V_z\bm{Z}\otimes\bm{Z}+V_x\bm{X}\otimes\bm{X}.
\]
\end{frame}

\begin{frame}[plain,fragile]
\frametitle{How do we perform measurements?}

The above tensor products have to be rewritten in terms of specific
transformations so that we can perform the measurements in the basis of
the Pauli-$\bm{Z}$ matrices. As we discussed earlier, we need to find
a transformation of the form
\[
\mathcal{P}=\bm{U}^{\dagger}\bm{M}\bm{U},
\]
where $\mathcal{P}$ represents some combination of the Pauli matrices and
the identity matrix, $\bm{U}$ is a unitary matrix and $\bm{M}$
represents the gate/matrix which performs the measurements, often
represented by a Pauli-$\bm{Z}$ gate/matrix.
\end{frame}

\begin{frame}[plain,fragile]
\frametitle{Rewriting our strings of Pauli matrices}

As discussed last week and reviewed above, to perform a measurement,
it is often useful to rewrite the string of Pauli matrices in a
specific order or to simplify the expression.

Consider a string of Pauli matrices of the form:
\[
P = \sigma_{i_1} \otimes \sigma_{i_2} \otimes \cdots \otimes \sigma_{i_n},
\]
where each $\sigma_{i_k}$ is one of the Pauli matrices $\bm{X}$, $\bm{Y}$, $\bm{Z}$, or the identity matrix $\bm{I}$.
\end{frame}

\begin{frame}[plain,fragile]
\frametitle{Rewriting the string of matrices}

To rewrite this string for measurement purposes, follow these steps:
\begin{block}{Commute and reorder: }
Use the commutation relations of the Pauli matrices to reorder the string. The commutation relations are:
\[
[\sigma_i, \sigma_j] = 2i \epsilon_{ijk} \sigma_k,
\]
where $\epsilon_{ijk}$ is the Levi-Civita symbol. This allows you to move certain Pauli matrices to the left or right in the string.
\end{block}
\begin{block}{Simplify using identities: }
Use the fact that $\sigma_i^2 = I$ and the above mentioned commutation relation.
\end{block}
\end{frame}

\begin{frame}[plain,fragile]
\frametitle{More rewriting tricks}

\begin{block}{Group terms: }
Group terms that are easier to measure together. For example, if you have a term like $\bm{X} \otimes \bm{X}$, you can measure both qubits in the $\bm{X}$-basis simultaneously.
\end{block}
\begin{block}{Diagonalize if necessary: }
If the final expression is not diagonal, you may need to apply a unitary transformation to diagonalize it before measurement. For example,to measure $\bm{X}$, you can apply the Hadamard gate $\bm{H}$ to transform it into $\bm{X}$:
\[
\bm{H} \bm{X} \bm{H} = \bm{Z}.
\]
\end{block}
\end{frame}

\begin{frame}[plain,fragile]
\frametitle{The $\bm{Z}\otimes \bm{I}$ term}

The explicit matrix is 
\[
\bm{Z}\otimes \bm{I} = \begin{bmatrix}
		1 & 0 & 0 & 0 \\
                0 & 1 & 0 & 0 \\
                0 & 0 & -1 & 0 \\
		0 & 0 & 0 & -1
	\end{bmatrix}.
\]

When we perform a measurement on the first qubit, we see that this
matrix gives us the correct eigenvalues for the first qubit (but not
for the second one). To see this multiply the above matrix with our
computational basis states, that is the states $\vert 00\rangle =\vert 0\rangle
\times \vert 0\rangle$, $\vert 01 \rangle$, $\vert 10\rangle$ and
$\vert 11\rangle$.

This process is referred to in the language of Pauli measurements as
\textbf{measuring Pauli-Z} and is entirely equivalent to performing a
computational basis measurement.
\end{frame}

\begin{frame}[plain,fragile]
\frametitle{The specific eigenvalues}

Multiplying with the $\vert 00\rangle$ state we get
\[
\bm{Z}\otimes \bm{I} = \begin{bmatrix}
		1 & 0 & 0 & 0 \\
                0 & 1 & 0 & 0 \\
                0 & 0 & -1 & 0 \\
		0 & 0 & 0 & -1
	\end{bmatrix}\begin{bmatrix} 1\\ 0 \\ 0 \\ 0\end{bmatrix}=\begin{bmatrix} 1\\ 0 \\ 0 \\ 0\end{bmatrix},
\]
and 
\[
\bm{Z}\otimes \bm{I} = \begin{bmatrix}
		1 & 0 & 0 & 0 \\
                0 & 1 & 0 & 0 \\
                0 & 0 & -1 & 0 \\
		0 & 0 & 0 & -1
	\end{bmatrix}\begin{bmatrix} 0\\ 0 \\ 1 \\ 0\end{bmatrix}=-1\begin{bmatrix} 0\\ 0 \\ 1 \\ 0\end{bmatrix}.
\]
We don't get the correct eigenvalues if we perform the measurement on the second qubit!
\end{frame}

\begin{frame}[plain,fragile]
\frametitle{The $\bm{I}\otimes\bm{Z}$ term}

Our strategy is to rewrite all the strings of Pauli operators via
specific unitary transformations in order to obtain a final operator
$\bm{P}=\bm{Z}_1\otimes\bm{I}_2\otimes\bm{I}_3\otimes\cdots\otimes \bm{I}_N$, where the subscripts indicate
the various qubits. We can then perform the measurement on the first qubit only.

We can rewrite $\bm{I}\otimes \bm{Z}$ via the SWAP gate
\[
\text{SWAP} = \begin{bmatrix}
		1 & 0 & 0 & 0 \\
                0 & 0 & 1 & 0 \\
                0 & 1 & 0 & 0 \\
		0 & 0 & 0 & 1
	\end{bmatrix}.
\]
\end{frame}

\begin{frame}[plain,fragile]
\frametitle{More on the $\bm{I}\otimes\bm{Z}$ term}

We have then that
\[
\bm{P}=\text{SWAP}^{\dagger}(\bm{I}\otimes\bm{Z})\text{SWAP}=\bm{Z}\otimes\bm{I}.
\]

We note that the original $\bm{I}\otimes \bm{Z}$ does not give the
correct eigenvalues when measured on the first qubit. Try this as an exercise.
\end{frame}

\begin{frame}[plain,fragile]
\frametitle{The $\bm{Z}\otimes \bm{Z}$ term}

Thus the tensor products of two Pauli-$\bm{Z}$ operators forms a matrix
composed of two spaces consisting of $+1$ and $-1$ as eigenvalues.
As with the single-qubit case, both constitute a half-space, meaning that half of
the accessible vector space belongs to the eigenspace with eigenvalue $+1$ and the
remaining half to the eigenspace with eigenvalue $-1$.

This term gives the correct eigenvalue when operating on the first
qubit. In principle thus we don't need to rewrite string of operators.
However, as discussed below as well, let us rewrite it via a unitary transformation in
order to have $\bm{P}=\bm{Z}\otimes\bm{I}$.  To do so, consider the
transformation

\[
\bm{P}= CX_{10}^{\dagger}(\bm{Z}\otimes\bm{Z})CX_{10},
\]
where we have 
\[
	\text{CX}_{10} = \begin{bmatrix}
		1 & 0 & 0 & 0 \\
		0 & 0 & 0 & 1 \\
		0 & 0 & 1 & 0 \\
		0 & 1 & 0 & 0
	\end{bmatrix}.
\]
\end{frame}

\begin{frame}[plain,fragile]
\frametitle{Performing the transformation}

Performing the operation gives
\[
\bm{P}= CX_{10}^{\dagger}(\bm{Z}\otimes\bm{Z})CX_{10}=\begin{bmatrix}
		1 & 0 & 0 & 0 \\
		0 & 1 & 0 & 0 \\
		0 & 0 & -1 & 0 \\
		0 & 0 & 0 & -1
	\end{bmatrix}=\bm{Z}\otimes\bm{I},
\]

which allows us to transform the original tensor product
$\bm{Z}\otimes \bm{Z}$ into a matrix where we perform the measurement
in the basis of the first qubit.

To see this, act with $\bm{P}$ on the states $\vert 00\rangle =\vert
0\rangle \times \vert 0\rangle$, $\vert 01 \rangle$, $\vert 10\rangle$
and $\vert 11\rangle$.
\end{frame}

\begin{frame}[plain,fragile]
\frametitle{Transformations}

Any unitary transformation of such matrices also describes two
half-spaces labeled with eigenvalues. For example
\[
\bm{X}\otimes\bm{X}=\bm{H}\otimes\bm{H}(\bm{Z}\otimes\bm{Z})\bm{H}\otimes\bm{H},, 
\]

follows from from the identity that $\bm{Z}=\bm{HXH}$.  Similar to the
one-qubit case, all two-qubit Pauli-measurements can be written in
terms of unitary transformations
$\bm{U}^{\dagger}(\bm{Z}\otimes\bm{I})\bm{U}$ with $\bm{U}$ being
$4\times 4$ unitary matrices.
\end{frame}

\begin{frame}[plain,fragile]
\frametitle{More terms}

Let us look at the $\bm{X}\otimes \bm{X}$ term in the simpler two-qubit Hamiltonian.

This term can be rewritten as

\[
\bm{P}= (CX_{10}(\bm{H}\otimes\bm{H}))^{\dagger}(\bm{X}\otimes\bm{X})(CX_{10}(\bm{H}\otimes\bm{H})),
\]
which results in 
\[
	\begin{bmatrix}
		1 & 0 & 0 & 0 \\
		0 & 1 & 0 & 0 \\
		0 & 0 & -1 & 0 \\
		0 & 0 & 0 & -1
	\end{bmatrix}.
\]
We recognize this result as our $\bm{Z}\otimes\bm{I}$ tensor product.
\end{frame}

\begin{frame}[plain,fragile]
\frametitle{Explicit expressions}

In order to perform our measurements for our simple two-qubit
Hamiltonian we need the following unitary transformations $\bm{U}$
\begin{align*}
\bm{Z}\otimes\bm{I}\hspace{1cm} & \bm{U}=\bm{I}\otimes\bm{I}\\
\bm{I}\otimes\bm{Z}\hspace{1cm} & \bm{U}=\text{SWAP}\\
\bm{Z}\otimes\bm{Z}\hspace{1cm} & \bm{U}=CX_{10}\\
\bm{X}\otimes\bm{X}\hspace{1cm} & \bm{U}=CX_{10}(\bm{H}\otimes\bm{H})\\
\end{align*}

These are the gates that are relevant for the simpler two-qubit Hamiltonian of project 1.
\end{frame}

\begin{frame}[plain,fragile]
\frametitle{More complete list and derivations of expressions for strings of operators}

For a two qubit system we list here the possible transformations 
\begin{align*}
\bm{Y}\otimes\bm{I}\hspace{1cm} & \bm{U}=\bm{H}\bm{S}^{\dagger}\otimes\bm{I}\\
\bm{I}\otimes\bm{X}\hspace{1cm} & \bm{U}=(\bm{H}\otimes\bm{I})\text{SWAP}\\
\bm{I}\otimes\bm{X}\hspace{1cm} & \bm{U}=(\bm{H}\otimes\bm{I})\text{SWAP}\\
\bm{I}\otimes\bm{X}\hspace{1cm} & \bm{U}=(\bm{H}\bm{S}^{\dagger}\otimes\bm{I})\text{SWAP}\\
\bm{X}\otimes\bm{Z}\hspace{1cm} & \bm{U}=CX_{10}(\bm{H}\otimes\bm{I})\\
\bm{Y}\otimes\bm{Z}\hspace{1cm} & \bm{U}=CX_{10}(\bm{H}\bm{S}^{\dagger}\otimes\bm{I})\\
\bm{Z}\otimes\bm{Y}\hspace{1cm} & \bm{U}=CX_{10}(\bm{I}\otimes\bm{H}\bm{S}^{\dagger})\\
\bm{Y}\otimes\bm{X}\hspace{1cm} & \bm{U}=CX_{10}(\bm{H}\bm{S}^{\dagger}\otimes\bm{H})\\
\bm{X}\otimes\bm{Y}\hspace{1cm} & \bm{U}=CX_{10}(\bm{H}\otimes\bm{H}\bm{S}^{\dagger})\\
\bm{Y}\otimes\bm{Y}\hspace{1cm} & \bm{U}=CX_{10}(\bm{H}\bm{S}^{\dagger}\otimes\bm{H}\bm{S}^{\dagger})\\
\end{align*}
\end{frame}

\begin{frame}[plain,fragile]
\frametitle{Additional remarks}

Here, the CNOT (CX10) operation appears for the following reason. Each Pauli
measurement that does not include the identity matrix is equivalent up to a
unitary to $\bm{Z}\otimes\bm{Z}$. The eigenvalues $\bm{Z}\otimes\bm{Z}$ of only depend on
the parity of the qubits that comprise each computational basis
vector, and the controlled-not operations serve to compute this parity
and store it in the first bit. Then once the first bit is measured,
you can recover the identity of the resultant half-space, which is
equivalent to measuring the Pauli operator.
\end{frame}

\begin{frame}[plain,fragile]
\frametitle{Reducing space}

Also, while it can be tempting to assume that measuring $\bm{Z}\otimes\bm{Z}$ is the same as
sequentially measuring $\bm{Z}\otimes\bm{I}$ and then $\bm{I}\otimes\bm{Z}$, this assumption would be
false. The reason is that measuring $\bm{Z}\otimes\bm{Z}$ projects the quantum state into
either the $+1$ or $-1$ eigenstate of these operators. Measuring $\bm{Z}\otimes\bm{I}$ and then
$\bm{I}\otimes\bm{Z}$
projects the quantum state vector first onto a half space of $\bm{Z}\otimes\bm{I}$ and
then onto a half space of $\bm{I}\otimes\bm{Z}$. As there are four computational basis
vectors, performing both measurements reduces the state to a
quarter-space and hence reduces it to a single computational basis
vector.
\end{frame}

\begin{frame}[plain,fragile]
\frametitle{Lipkin model}

We will study a schematic model (the Lipkin model, see Nuclear
Physics \textbf{62} (1965) 188), for the interaction among  $2$ and more 
fermions that can occupy two different energy levels.

In project 1 we consider a two-fermion case and a four-fermion case.

For four fermions, the case we consider in the examples  here, each levels has
degeneration $d=4$, leading to different total spin values.  The two
levels have quantum numbers $\sigma=\pm 1$, with the upper level
having $2\sigma=+1$ and energy $\varepsilon_{1}= \varepsilon/2$. The
lower level has $2\sigma=-1$ and energy
$\varepsilon_{2}=-\varepsilon/2$. That is, the lowest single-particle
level has negative spin projection (or spin down), while the upper
level has spin up.  In addition, the substates of each level are
characterized by the quantum numbers $p=1,2,3,4$.
\end{frame}

\begin{frame}[plain,fragile]
\frametitle{Four fermion case}

We define the single-particle states (for the four fermion case which we will work on here)
\[
\vert u_{\sigma =-1,p}\rangle=a_{-p}^{\dagger}\vert 0\rangle
\hspace{1cm}
\vert u_{\sigma =1,p}\rangle=a_{+p}^{\dagger}\vert 0\rangle.
\]
The single-particle states span an orthonormal basis.
\end{frame}

\begin{frame}[plain,fragile]
\frametitle{Hamiltonian}

The Hamiltonian of the system is given by

\[
\begin{array}{ll}
\hat{H}=&\hat{H}_{0}+\hat{H}_{1}+\hat{H}_{2}\\
&\\
\hat{H}_{0}=&\frac{1}{2}\varepsilon\sum_{\sigma ,p}\sigma
a_{\sigma,p}^{\dagger}a_{\sigma ,p}\\
&\\
\hat{H}_{1}=&\frac{1}{2}V\sum_{\sigma ,p,p'}
a_{\sigma,p}^{\dagger}a_{\sigma ,p'}^{\dagger}
a_{-\sigma ,p'}a_{-\sigma ,p}\\
&\\
\hat{H}_{2}=&\frac{1}{2}W\sum_{\sigma ,p,p'}
a_{\sigma,p}^{\dagger}a_{-\sigma ,p'}^{\dagger}
a_{\sigma ,p'}a_{-\sigma ,p}\\
&\\
\end{array}
\]
where $V$ and $W$ are constants. The operator 
$H_{1}$ can move pairs of fermions
while $H_{2}$ is a spin-exchange term. The latter
moves a pair of fermions from a state $(p\sigma ,p' -\sigma)$ to a state
$(p-\sigma ,p'\sigma)$.
\end{frame}

\begin{frame}[plain,fragile]
\frametitle{Quasispin operators}

We are going to rewrite the above Hamiltonian in terms of so-called  quasispin operators
\[
\begin{array}{ll}
\hat{J}_{+}=&\sum_{p}
a_{p+}^{\dagger}a_{p-}\\
&\\
\hat{J}_{-}=&\sum_{p}
a_{p-}^{\dagger}a_{p+}\\
&\\
\hat{J}_{z}=&\frac{1}{2}\sum_{p\sigma}\sigma
a_{p\sigma}^{\dagger}a_{p\sigma}\\
&\\
\hat{J}^{2}=&J_{+}J_{-}+J_{z}^{2}-J_{z}\\
&\\
\end{array}
\]
These operators obey the commutation relations for angular momentum.
\end{frame}

\begin{frame}[plain,fragile]
\frametitle{Including the number operator}

We can in turn express $\hat{H}$ in terms of the above quasispin operators and the number operator
\[
\hat{N}=\sum_{p\sigma}
a_{p\sigma}^{\dagger}a_{p\sigma}.
\]
\end{frame}

\begin{frame}[plain,fragile]
\frametitle{Final  Hamiltonian}

Using the quasispin operators we can rewrite the Hamiltonian as
\begin{equation}
H = \varepsilon J_z+\frac{1}{2} V \left( J_+^2 + J_-^2 \right)+\frac{1}{2} W \left( -N + J_+ J_- + J_- J_+ \right).
\end{equation}
\end{frame}

\begin{frame}[plain,fragile]
\frametitle{Final Hamiltonian matrix}

These five states can in turn be used as computational basis states in
order to define the Hamiltonian matrix to be diagonalized.
The matrix elements are given by $\langle J,J_z \vert H \vert J',J_z' \rangle$.
The
Hamiltonian is hermitian and we obtain after all this labor of ours

\begin{equation}
H_{J = 2} =
\begin{bmatrix}
-2\varepsilon & 0 & \sqrt{6}V & 0 & 0 \\
0 & -\varepsilon + 3W & 0 & 3V & 0 \\
\sqrt{6}V & 0 & 4W & 0 & \sqrt{6}V \\
0 & 3V & 0 & \varepsilon + 3W & 0 \\
0 & 0 & \sqrt{6}V & 0 & 2\varepsilon
\end{bmatrix}
\label{eq:HJ=2}
\end{equation}
\end{frame}

\begin{frame}[plain,fragile]
\frametitle{Quantum computing and solving  the eigenvalue problem for the Lipkin model}

We turn now to a simpler variant of the Lipkin model without the $W$-term and a total spin of $J=1$ only as maximum value of the spin.
This corresponds to a system with $N=2$ particles (fermions in our case).
Our Hamiltonian is given by the quasispin operators (see below) 
\[
     \hat{H} = \epsilon\hat{J}_z -\frac{1}{2}V(\hat{J}_+\hat{J}_++\hat{J}_-\hat{J}_-).
\]

As discussed earlier
the quasispin operators act like lowering and raising angular momentum
operators.

With these properties we can calculate the Hamiltonian
matrix for the Lipkin model by computing the various matrix elements
\begin{equation}
\langle JJ_z|H|JJ_z'\rangle,
\end{equation}
where the non-zero elements are given by
\[
\begin{split}
\langle JJ_z|H|JJ_z'\rangle & = \epsilon J_z\\
\langle JJ_z|H|JJ_z'\pm 2\rangle & = \langle JJ_z\pm 2|H|JJ_z'\rangle \\ &= -\frac{1}{2}VC,
\end{split}
\]
where $C$ is the Clebsch-Gordan coefficients (from the raising and lowering operators) one gets when
$J_{\pm}^2$ operates on the state $|JJ_z\rangle$.  Using the above
definitions we can calculate the exact solution to the Lipkin model.

With the $V$-interaction terms, we obtain the following Hamiltonian matrix
\begin{equation}
\begin{pmatrix}-\epsilon & 0 & -V\\
 0&0&0\\
 -V&0&\epsilon
\end{pmatrix}
\end{equation}
\end{frame}

\begin{frame}[plain,fragile]
\frametitle{Plans for next week}

\begin{enumerate}
\item We will focus mainly on setting up the quantum circuit for the Lipkin model with $J=1$ and $J=2$, rewriting the Lipkin model in terms of Pauli matrices 

\item Thereafter, if we get time we start with quantum Fourier transforms and discussions of Quantum Phase estimation algorithm
\end{enumerate}

\noindent
\end{frame}

\end{document}
